\documentclass[10pt]{article}
\usepackage[margin=1in]{geometry}
\usepackage[T1]{fontenc}
\usepackage{lmodern,microtype,amsmath,amssymb,amsthm}
\theoremstyle{definition}
\newtheorem{assumption}{Assumption}
\begin{document}
\section*{Constrained Generalized Monotone Variational Inequalities}
We consider a variational inequality for a continuous monotone operator $F:E\to E$ on a convex set $X$ in a real inner product space $E$. Operator variation satisfies a generalized Lipschitz bound controlled by its norm along each segment. A feasible initial point $x^0$ has bounded distance to a solution, with bounded operator norms at both points. Two separately counted oracles return exact operator values and metric projections onto $X$. For $\varepsilon>0$, the output $\bar{x}$ must belong to $X$ and admit $\xi\in E$ satisfying $\langle\xi,w-\bar{x}\rangle\leq0$ for every $w\in X$ and $\|F(\bar{x})+\xi\|\leq\varepsilon$, almost surely under the joint law of algorithm and oracle randomness.

The underlying variational inequality is to find $x^\star\in X$ such that
\[
\langle F(x^\star),w-x^\star\rangle\geq0
\qquad\text{for every }w\in X.
\]
The norm on $E$ is induced by its inner product; $E$ need not be finite dimensional or complete. The parameters satisfy
\[
L_0>0,\qquad L_1\geq0,\qquad R>0,\qquad
G\geq0,\qquad G_0\geq0,\qquad\varepsilon>0.
\]
The admissible instances and oracle access satisfy the following conditions.

\begin{assumption}[Convex feasibility and monotonicity]
The set $X\subseteq E$ is convex, the operator $F:E\to E$ is continuous, and
\[
\langle F(x)-F(y),x-y\rangle\geq0
\qquad\text{for every }x,y\in E.
\]
\end{assumption}

\begin{assumption}[Generalized Lipschitz continuity]
For every $x,y\in E$ and every $M\in\mathbb{R}$, the segment bound
\[
\|F(x+\theta(y-x))\|\leq M
\qquad\text{for every }\theta\in[0,1]
\]
implies
\[
\|F(x)-F(y)\|
\leq (L_0+L_1M)\|x-y\|.
\]
\end{assumption}

\begin{assumption}[Solution and initialization bounds]
The initial point $x^0$ belongs to $X$, and there exists a solution $x^\star\in X$ of the variational inequality such that
\[
\begin{aligned}
\|x^0-x^\star\|&\leq R,\\
\|F(x^0)\|&\leq G_0,\\
\|F(x^\star)\|&\leq G.
\end{aligned}
\]
\end{assumption}

\begin{assumption}[Exact operator and projection access]
Oracle $0$ returns $F(x)\in E$ at any query $x\in E$. Oracle $1$ returns $\Pi_X(z)\in E$ at any query $z\in E$, where a map $\Pi_X:E\to E$ exists and satisfies
\[
\begin{aligned}
\Pi_X(z)&\in X,\\
\langle z-\Pi_X(z),w-\Pi_X(z)\rangle&\leq0
\qquad\text{for every }w\in X
\end{aligned}
\]
for every $z\in E$. Thus $\Pi_X(z)$ is the metric projection of $z$ onto $X$. Both replies are exact at every call index and for every realization of oracle randomness. Each query costs one call to the queried oracle, with calls counted separately. Internal algorithm randomness is independent of oracle randomness.
\end{assumption}

The algorithm is selected using only the public inputs $(E,L_0,L_1,x^0)$ and accesses $F$ and $X$ through these oracles. Queries may be adaptive and may lie anywhere in $E$, including outside $X$. The quantities $R$, $G$, $G_0$, and $\varepsilon$ are not supplied to the algorithm. Externally specified budgets $K_0,K_1\in\mathbb{N}$ may depend on the full parameter tuple $(E,L_0,L_1,x^0,R,G,G_0,\varepsilon)$ and limit calls to oracles $0$ and $1$, respectively.

The algorithm returns $\bar{x}\in E$. Writing $\mathbb{P}$ for the joint law of algorithm and oracle randomness, the required accuracy criterion is
\[
\begin{gathered}
\bar{x}\in X,\qquad \exists\,\xi\in E:\\
\langle\xi,w-\bar{x}\rangle\leq0
\quad\text{for every }w\in X,\\
\|F(\bar{x})+\xi\|\leq\varepsilon,
\end{gathered}
\qquad\mathbb{P}\text{-almost surely}.
\]
The vector $\xi$ is a witness for the accuracy criterion; the required output is the point $\bar{x}$.
\end{document}
